\section{Negative Euler certificates on partial quadrilateral masks}
\label{sec:euler-branch-tree}

One-vertex preprocessing removes the Cartesian anchor-product obstruction,
but a flat negative transcript still names all $3^t$ maximal compatible
sectors.  The branch-tree interface permits one checked relaxation to
exclude many such sectors simultaneously.  A negative certificate can stop
at a partial quadrilateral assignment; it need not expand every remaining
choice merely to satisfy the global verifier.

This is an implementation of established relaxation and branching ideas,
not a new general branching theorem.  The preserved standard-coordinate
trees, mandatory propagation and deciding-set bounds of report~67 are
discussed in Section~\ref{sec:completion-theorems}; report~73 additionally
studies active-support certificates and conflict-rank branching in
Section~\ref{sec:incoming-contexts}.  The interface here
uses the maintained Q-potential model and its envelope certificates, with
one independently reconstructed source and an explicitly checked coverage
tree.  The distinction between formal relaxation points and admissible
normal surfaces is essential throughout.

\subsection{A negative statement remains sound without quad compatibility}

Let $J_t\subseteq\{0,1,2\}$ be the allowed types in tetrahedron $t$.
An empty set forces all its quadrilateral coordinates to zero; a singleton
is an ordinary compatible sector restriction.  A larger set permits
incompatible coordinates in the relaxation.  Write $J$ for the sorted list
of all allowed pairs and $k=|J|\le3t$.

Triangle matching and vertex-link peeling remain linear constructions
before compatibility is imposed.  The source-derived cycle equations give
\[
 Q_J=\{q\in\mathbb R^k:q\ge0,\ C_Jq=0\}.
\]
For every $q\in Q_J$, triangle potentials and subtraction of their minimum
at each global vertex give a nonnegative canonical matching vector.  This
vector represents an embedded normal surface only when its quad coordinates
are compatible.  Nevertheless, the ordinary source Euler coefficients
define a linear functional on all matching vectors.  Evaluating that
functional on the canonical lift gives the formal extension
\[
 F_J(q)=L_J(q)-\sum_v w_v\min_{c\in G_v}\ell_{vc}(q)
       =\max_{\alpha\in\prod_vG_v}c_\alpha^{\mathsf T}q,
 \qquad w_v\in\{1,2\}.
\]
The identity uses only matching and the vertex-link coefficients.  No
quadrilateral compatibility hypothesis enters its algebraic derivation.

\begin{proposition}[Relaxation exclusion]
If $F_J(q)\le0$ for every $q\in Q_J$, no admissible canonical positive-Euler
normal surface has all its positive quadrilateral types in $J$.
\end{proposition}
\begin{proof}
Every such surface projects to a point of $Q_J$.  Its canonical triangle
coordinates agree with the source-derived canonical lift, and on this
admissible vector the formal Euler functional equals the surface's Euler
characteristic.  A positive characteristic would contradict the assumed
nonpositivity.  Thus a proof on the larger incompatible relaxation excludes
every compatible sub-sector it contains.
\end{proof}

The same homogeneous Farkas checks from
Section~\ref{sec:euler-lp-crushing} suffice.  For every anchor objective one
may supply $y_\alpha$ with
$C_J^{\mathsf T}y_\alpha\ge c_\alpha$.  Alternatively a single multiplier
may dominate the componentwise envelope of all those objectives.  The
checker reconstructs the matching model from the original face pairing.
Allowing several distinct types per tetrahedron changes the support parser,
not the negative inequality proof.

The dedicated \texttt{normal-euler-relaxation-v1} schema is negative-only.
It names the source digest and the sorted, distinct allowed pairs, then
contains either the envelope multipliers or the complete anchor-dual list.
Literal integer indices, exact rational arithmetic and complete anchor
coverage remain required.  The existing compatible-sector positive schema
retains its original restrictions.

\subsection{A fixed source model permits checked reuse of dual bounds}

Reconstructing and contracting the Q model separately at every mask is
unnecessary.  Let $U$ contain all $3t$ quadrilateral types, and construct once
from the source its cycle matrix $C$, Euler row $L$, triangle potentials
$\ell_{vc}$ and corner groups $G_v$.  For $J\subseteq U$, let $\iota_J$ insert
zeroes in the omitted coordinates.  Then
\[
 C\iota_Jq=C[:,J]q,\qquad
 F(\iota_Jq)=\max_{\alpha\in\prod_vG_v}c_\alpha[J]^{\mathsf T}q.
\]
Thus a node can restrict columns of one fixed matrix and one fixed family
of objectives.  It keeps every source cycle row, including rows whose
restriction becomes zero.  This preserves the interpretation of a dual
multiplier at every node.

For completeness, this fixed model describes the same canonical vectors
as a freshly reconstructed restricted model.  A full-model Q vector
supported on $J$ has a triangle lift precisely when its zero expansion
satisfies the source cycle equations.  Both constructions solve the same
original triangle matching equations with the same quadrilateral values.
Their triangle solutions differ only by vertex-link constants, so
subtracting the minimum at each global vertex gives identical canonical
coordinates and Euler value.  Some fixed potentials can coincide after
restriction; retaining all their original anchor choices merely repeats
equal objectives.  It does not introduce a new positivity obligation.
In particular, on a one-vertex source the full fixed family has at most
$4t$ choices, regardless of how many further identifications a node-specific
construction could discover.

Suppose a solved negative query gives a rational dual $y$.  Compute and
retain its full image $b=C^{\mathsf T}y$.  For any later mask $J$ and
objective $d$, this same multiplier proves nonpositivity whenever
\[
 b_j\ge d_j\quad\hbox{for every }j\in J.
\]
Indeed, for every nonnegative $q$ with $C[:,J]q=0$,
$d[J]^{\mathsf T}q\le b[J]^{\mathsf T}q
=y^{\mathsf T}C[:,J]q=0$.
No inequality outside $J$ is required.  If the objective is unchanged and
the mask shrinks, the inequalities are inherited.  Reuse from a different
branch, or for a different anchor objective, is equally sound only after
the current retained-column inequalities have been checked again.  A
previously successful query is not itself evidence for these new
inequalities.

The componentwise envelope is also fixed at the source:
\[
 e_j=L_j-\sum_v w_v\min_{c\in G_v}\ell_{vc,j}.
\]
It dominates every $c_{\alpha,j}$, so the same cache lookup can discharge an
envelope query by testing $b[J]\ge e[J]$.  Conversely an envelope dual may
later dominate an individual anchor.  The zero multiplier is a useful
permanent cache entry: it immediately handles any objective nonpositive on
the retained columns.  The implementation retains at most $256$ entries,
including this zero entry, across the search on one source.  Bounded
retention makes each lookup a bounded number of polynomial-size exact
comparisons rather than an ever-growing scan over the search history.
Eviction can change discovery work, but cannot change correctness.  A
crushing successor begins a new source model and a new cache; source rows
and multipliers cannot be transferred merely because their dimensions agree.

The \texttt{normal-euler-relaxation-v2} leaf format has the same fields as
version~1, but defines its multiplier coordinates and anchor enumeration
using the full source model.  Its independent checker reconstructs that
model at most once per tree, verifies the complete fixed anchor family
when required, and tests inequalities only on the leaf's allowed columns.
Version~1 retains its original mask-specific reconstruction for comparison.
The versions cannot be interchanged without checking their different row
and anchor conventions.  Both carry actual exact multipliers: the verifier
does not trust cache identifiers, cache history or a producer's claim that
a dual was reused successfully.  Consequently a tree may contain either
version without altering its coverage argument.

\subsection{Positive formal points guide discovery, not topology}

A point with $F_J(q)>0$ need not describe a surface: two positive quad types
in one tetrahedron already violate embeddedness.  Such a point may guide
the next branching choice, but it cannot be passed to crushing as a normal
witness.  The source coordinate validator and positive support extractor
remain the gate to that operation.  A compatible positive canonical point
can use their existing contracts; an incompatible point must be resolved
by further restriction.

If tetrahedron $a$ currently permits at least two types, split it into one
child for each $q\in J_a$, replacing $J_a$ by $\{q\}$ and retaining every
other mask entry.  Every admissible surface in the parent belongs to a
child: its positive quad type determines a child, or, if all three quad
coordinates there are zero, it belongs to every child.  Consequently a
separate zero-quad child is unnecessary.  The children may overlap on such
zero-coordinate faces; disjointness is not needed to prove exclusion.

A positive envelope objective alone need not even give positive $F_J$.
The actual formal canonical value must be checked as in the maintained
Euler producer; otherwise an unsuccessful relaxation merely leaves the
negative obligation unresolved.  A local pivot or node limit likewise
leaves an unfinished search, rather than producing a negative leaf.

\subsection{Implicit contexts make coverage independently checkable}

The independent checker is
\path{fast/fastunknot/normal_euler_tree_verify.py}.
The \texttt{normal-euler-branch-tree-v1} packet names its source and root
allowed pairs and stores a flat preorder sequence of nodes.  A split node
contains only its kind and the tetrahedron index.  Its children are implicit:
the verifier derives the singleton masks from the current context in
ascending allowed-type order.  A negative node contains one independently
checked relaxation certificate, whose source and allowed mask must equal
that context exactly.  This tree schema has no positive leaves or nested
matching-reduction wrappers.

An iterative stack initially contains the root context.  Each node consumes
the next pending context.  A legal split replaces it by the derived child
contexts in preorder; a negative leaf discharges it.  Verification succeeds
only if the node list and pending contexts are both exhausted together.
Missing children, extra nodes, a split at an already singleton mask,
reordered child masks, and proofs for a foreign source or smaller mask
therefore fail structurally.  Leaf inequalities are still checked separately;
correct tree syntax cannot replace them.

\begin{proposition}[Sound negative tree]
A verified tree excludes every admissible canonical positive-Euler normal
surface in its root mask.
\end{proposition}
\begin{proof}
A negative leaf has this property by relaxation exclusion.  If all children
of a split exclude their masks, admissible coverage by the singleton
children gives the same conclusion for the parent.  Induction from the
leaves to the root proves the assertion.  The preorder stack check ensures
that every required child participates in this induction exactly where its
derived context requires it.
\end{proof}

For global recognition the root must allow all three types in every
tetrahedron of the final authenticated source.  A valid tree on a partial
root mask is only a restricted statement.  The diagram-bound verifier must
check this full root and the preceding crushing chain before using the
absence of positive canonical surfaces to conclude knottedness.  The normal
disc existence theorem and the source correspondence remain those of
Section~\ref{sec:normal-crushing-search}; tree compression changes the
coverage evidence, not those topological premises.

\subsection{Complexity of compression and its remaining limit}

On each one-vertex source supplied by
Section~\ref{sec:vertex-reduction}, there is at most one active triangle
group.  Even for an incompatible mask, its number of corner classes is at
most $4t$.  Thus a complete node has at most $4t$ anchor objectives, or one
empty objective when there is no active group, with at most one additional
envelope call.  The number of allowed Q coordinates is at most $3t$.
The fixed source model is prepared once.  Forming each column restriction,
checking its bounded dual cache and checking each rational inequality take
polynomial bit work in the source and rational data involved.  A complete
tree is independently checkable in time polynomial in the source size and
the encoded certificate size.  Version~2 avoids reconstructing a model at
every leaf, but still records all required anchor multipliers; the cache is
a discovery optimization, not an implicit certificate-coverage rule.

Each split makes one previously ambiguous tetrahedron singleton, so depth
is at most $t$.  The general node bound is still
\[
 N\le 1+3+\cdots+3^t=\frac{3^{t+1}-1}{2}.
\]
With a polynomial exact LP policy, this would give
$N\operatorname{poly}(t)$ discovery on a one-vertex source.  The maintained
Bland implementation has no polynomial pivot bound, and no theorem here
bounds $N$ quasi-polynomially for all authenticated successors.

A negative leaf near the root can replace exponentially many explicit
sector packets.  That is a genuine certificate-size opportunity, and small
observed trees can be measured as such.  It is not a uniform tree-size
bound.  The preserved positive-relaxation and ambiguity obstructions also
remain relevant: one-vertex reduction removes their multi-vertex premise
where applicable, but does not force every remaining formal relaxation to
be nonpositive or every incompatible point to have a small deciding set.
